K_t^\Phi=K_t^\Phi\mathbf 1_{\{|x-y|<1\}}
        +K_t^\Phi\mathbf 1_{\{|x-y|\ge1\}},
\]
the near-diagonal majorant \eqref{eq:phi-cutoff-near-majorant}, the far-field
majorant \eqref{eq:phi-cutoff-far-majorant}, and their exchanged-variable
counterparts.  These are the exact majorants used above to prove
\[
\int_a^b\iint_{\RR^2}(1-\chi_R(x)\chi_R(y))K_t^\Phi(x,y)
\,\dd x\,\dd y\,\dd t\longrightarrow0
\]
for the first Bregman term, with the same argument after exchanging \(x\) and
\(y\).  Absolute convergence is the bound
\eqref{eq:doob-bregman-integrability-bound}; and the passage from smooth
convex entropies to the present range-local $C^{1,1}$ class uses the
approximants $\Phi_\varepsilon$ of
Lemma~\ref{lem:compact-range-phi-approximation}.  Their first derivatives have
a common Lipschitz constant on the compact quotient range, so the same
one-variable and Bregman majorants dominate both sides and dominated
convergence lets $\varepsilon\downarrow0$.  Convexity gives
$D_\Phi\ge0$, hence monotonicity.  Since $[a,b]$ was arbitrary, the proposition
follows.
\end{proof}

\begin{lemma}[Localized cutoff commutator for the Doob--Bregman identity]%
\label{lem:doob-bregman-cutoff-commutator}
Let the hypotheses and notation of
Lemma~\ref{lem:doob-bregman-domain-bounds} hold on
\([a,b]\Subset(0,\infty)\).  Put
\[
A_t(x):=\Phi'(u_t(x)),\qquad
G_t(x):=\Phi(u_t(x))-u_t(x)\Phi'(u_t(x)),
\]
and, for \(x\ne y\),
\begin{align*}
\mathcal B_t^\Phi(x,y)
&:=
g_t(x)D_\Phi(u_t(x),u_t(y))
+g_t(y)D_\Phi(u_t(y),u_t(x)),\\
\mathcal M_t^\Phi(x,y)
&:=
\bigl(h_t(x)-h_t(y)\bigr)A_t(y)
+\bigl(g_t(x)-g_t(y)\bigr)G_t(y).
\end{align*}
Let \(\chi_R(x)=\chi(x/R)\), where
\(\chi\in C_c^\infty(\RR)\), \(0\le\chi\le1\), and
\(\chi\equiv1\) on \([-1,1]\).  Define
\[
\mathcal C_R^\Phi(t,x,y)
:=
\chi_R(x)(1-\chi_R(y))\mathcal B_t^\Phi(x,y)
+(\chi_R(x)-\chi_R(y))\mathcal M_t^\Phi(x,y).
\]
Then the cutoff error in the localized entropy balance is exactly
\begin{equation}\label{eq:phi-cutoff-commutator-error-formula}
e_R(t)=
\frac{1}{2\pi}\iint_{\RR^2}
\frac{\mathcal C_R^\Phi(t,x,y)}{(x-y)^2}\dd x\dd y ,
\end{equation}
with the absolutely convergent near/far splitting
\begin{align}
e_R^{<}(t)
&:=
\frac{1}{2\pi}\iint_{|x-y|<1}
\frac{\mathcal C_R^\Phi(t,x,y)}{(x-y)^2}\dd x\dd y,
\label{eq:phi-cutoff-commutator-near}\\
e_R^{>}(t)
&:=
\frac{1}{2\pi}\iint_{|x-y|\ge1}
\frac{\mathcal C_R^\Phi(t,x,y)}{(x-y)^2}\dd x\dd y.
\label{eq:phi-cutoff-commutator-far}
\end{align}
Moreover
\begin{equation}\label{eq:phi-cutoff-commutator-l1-vanishing}
\int_a^b\bigl(|e_R^{<}(t)|+|e_R^{>}(t)|\bigr)\dd t
\le
\frac{1}{2\pi}\int_a^b\iint_{\RR^2}
\frac{|\mathcal C_R^\Phi(t,x,y)|}{|x-y|^2}\dd x\dd y\dd t
\longrightarrow0 .
\end{equation}
\end{lemma}

\begin{proof}
Write \(u_x=u_t(x)\), \(u_y=u_t(y)\), \(g_x=g_t(x)\),
\(g_y=g_t(y)\), \(h_x=g_xu_x\), \(h_y=g_yu_y\),
\(A_x=A_t(x)\), \(A_y=A_t(y)\), \(G_x=G_t(x)\), and \(G_y=G_t(y)\).
The cutoff bilinear computation for
\(\frac{\dd}{\dd t}H_R(t)\) produces the numerator
\[
\mathcal S_R(t,x,y)
:=
(h_x-h_y)\bigl(\chi_R(x)A_x-\chi_R(y)A_y\bigr)
+(g_x-g_y)\bigl(\chi_R(x)G_x-\chi_R(y)G_y\bigr).
\]
The desired product-cutoff numerator is
\[
\chi_R(x)\chi_R(y)
\Bigl[(h_x-h_y)(A_x-A_y)+(g_x-g_y)(G_x-G_y)\Bigr].
\]
By Lemma~\ref{lem:doob-bregman-algebra}, the bracket is
\(\mathcal B_t^\Phi(x,y)\).  Subtracting the product-cutoff numerator from
\(\mathcal S_R\) gives
\[
\begin{split}
&\mathcal S_R(t,x,y)-\chi_R(x)\chi_R(y)\mathcal B_t^\Phi(x,y)\\
&\qquad=
\chi_R(x)(1-\chi_R(y))\mathcal B_t^\Phi(x,y)
+(\chi_R(x)-\chi_R(y))
\bigl[(h_x-h_y)A_y+(g_x-g_y)G_y\bigr],
\end{split}
\]
which is exactly \(\mathcal C_R^\Phi(t,x,y)\).  This proves the formula for
the error once the bilinear identity is applied to the localized tests
\(\chi_R A_t\) and \(\chi_R G_t\).

It remains to justify the absolute convergence and the
\(L^1([a,b])\)-vanishing.  On the compact quotient range \(J_{a,b}\), the
functions \(\Phi\), \(\Phi'\), \(u_t\), \(A_t\), and \(G_t\) are uniformly
bounded, and \(A_t\), \(G_t\) are uniformly Lipschitz in \(x\).  The Bregman
part satisfies, by Lemma~\ref{lem:doob-bregman-domain-bounds},
\begin{align}
0\le \frac{\mathcal B_t^\Phi(x,y)}{|x-y|^2}
&\le C_{\Phi,a,b}(g_t(x)+g_t(y))\mathbf 1_{\{|x-y|<1\}},
\label{eq:phi-commutator-near-bregman-bound}\\
0\le \frac{\mathcal B_t^\Phi(x,y)}{|x-y|^2}
&\le C_{\Phi,a,b}(g_t(x)+g_t(y))|x-y|^{-2}
\mathbf 1_{\{|x-y|\ge1\}} .
\label{eq:phi-commutator-far-bregman-bound}
\end{align}
Both right-hand sides are integrable on
\([a,b]\times\RR^2\).

For the \(\mathcal M_t^\Phi\) part near the diagonal, the uniform boundedness
of \(A_t\) and \(G_t\) gives
\[
|\mathcal M_t^\Phi(x,y)|
\le C_{\Phi,a,b}\bigl(|h_t(x)-h_t(y)|+|g_t(x)-g_t(y)|\bigr).
\]
Using
\[
p_t(x)-p_t(y)
=(x-y)\int_0^1\partial_xp_t(y+s(x-y))\dd s,
\qquad p_t\in\{h_t,g_t\},
\]
and
\(|\chi_R(x)-\chi_R(y)|
\le \|\chi'\|_\infty |x-y|/R\le \|\chi'\|_\infty |x-y|\), we obtain
\begin{equation}\label{eq:phi-commutator-near-M-bound}
\frac{|(\chi_R(x)-\chi_R(y))\mathcal M_t^\Phi(x,y)|}{|x-y|^2}
\mathbf 1_{\{|x-y|<1\}}
\le
C_{\Phi,a,b,\chi}\int_0^1
\sum_{p\in\{h,g\}}|\partial_xp_t(y+s(x-y))|\dd s\,
\mathbf 1_{\{|x-y|<1\}} .
\end{equation}
The right-hand side is integrable on \([a,b]\times\RR^2\), because
Lemma~\ref{lem:cauchy-doob-cutoff-calculus} gives uniform \(L^1\)-bounds for
\(\partial_xh_t\) and \(\partial_xg_t\), and integration in the difference
variable contributes the finite factor \(\int_{|z|<1}\dd z=2\).
For each fixed \(t,x,y\), the factor
\(\chi_R(x)(1-\chi_R(y))\) and the factor
\(\chi_R(x)-\chi_R(y)\) tend to zero as \(R\to\infty\).  Dominated
convergence with
\eqref{eq:phi-commutator-near-bregman-bound} and
\eqref{eq:phi-commutator-near-M-bound} therefore gives
\[
\int_a^b\iint_{|x-y|<1}
\frac{|\mathcal C_R^\Phi(t,x,y)|}{|x-y|^2}
\dd x\dd y\dd t\longrightarrow0 .
\]

For the far field, boundedness of \(A_t\) and \(G_t\) gives
\[
|\mathcal M_t^\Phi(x,y)|
\le C_{\Phi,a,b}\bigl(h_t(x)+h_t(y)+g_t(x)+g_t(y)\bigr).
\]
Together with \(|\chi_R(x)-\chi_R(y)|\le1\) and
\eqref{eq:phi-commutator-far-bregman-bound}, this yields the integrable
majorant
\begin{equation}\label{eq:phi-commutator-far-M-bound}
\frac{|\mathcal C_R^\Phi(t,x,y)|}{|x-y|^2}
\mathbf 1_{\{|x-y|\ge1\}}
\le
C_{\Phi,a,b,\chi}
\frac{h_t(x)+h_t(y)+g_t(x)+g_t(y)}{|x-y|^2}
\mathbf 1_{\{|x-y|\ge1\}} .
\end{equation}
Indeed \(h_t\) and \(g_t\) are probability densities, so the right-hand side
has \([a,b]\times\RR^2\)-integral bounded by a constant times
\((b-a)\int_{|z|\ge1}|z|^{-2}\dd z\).  Pointwise again
\(\chi_R(x)(1-\chi_R(y))\to0\) and
\(\chi_R(x)-\chi_R(y)\to0\), so dominated convergence gives
\[
\int_a^b\iint_{|x-y|\ge1}
\frac{|\mathcal C_R^\Phi(t,x,y)|}{|x-y|^2}
\dd x\dd y\dd t\longrightarrow0 .
\]
Combining the near and far estimates proves
\eqref{eq:phi-cutoff-commutator-l1-vanishing}.
\end{proof}
